\section{Conclusion}
\label{sec:conclusion}

For offline, subject-independent gesture recognition with a two-ring monopolar forearm montage in able-bodied participants, four of 16 electrodes could be removed while remaining non-inferior to the full array within a prespecified 0.02 Macro-F1 engineering tolerance, but five could not. Cross-fitted attribution produced rankings that outperformed the mean of random subsets and whose faithfulness was supported by deletion with retraining, whereas the stability penalty produced no measurable improvement.

The study does not establish clinical performance, a reduction of physical bipolar hardware, an advantage of attribution over simple evenly spaced layouts, real-time performance, muscle-level localisation, or generalisation beyond these cohorts and montages. Validation in target clinical populations, such as people with upper-limb loss, and in online, causal control settings is required before deployment claims can be made.
